
%  The lines beginning with "%" are comments. 
%  Theses comments explain how to use this template for
%  preparing the Seminar report. 
%
%  The first non-comment line below contain the matter that
%  go into the preamble of the input file. These are collected in a 
%  file by name "preamble.tex". This file can be seen in the 
%  folder "SeminarReportTemplate". 
%
%  Do not make any changes in the first non-comment line below.
%
\documentclass[11pt, twoside]{book}
\let\cleardoublepage\clearpage
\usepackage{graphicx, fancyhdr, amsmath,times, url, makeidx}

\setlength{\oddsidemargin}{3cm}
\newcommand{\VAtitle}[1]{\def\vtitle{Bitcoin: A Cryptocurrency}}
\newcommand{\VAauthor}[1]{\def\vauthor{Savitha K.J.}}% Write full name without, repeat without, Mr. or Ms.
\newcommand{\VAadmissionyear}[1]{\def\vadmissionyear{2010}}
\newcommand{\VAregisternumber}[1]{\def\vregisternumber{VEAKECS049}}% Write University Examination Register Number.
\newcommand{\VAguide}[1]{\def\vguide{Ms. Geethu P C M.Tech, MISTE}}% Write full name with Mr. or Ms. or Dr. or Prof.
\newcommand{\VAhod}[1]{\def\vhod{Col MP Induchudan,  M.E., FITE}}% Write full name with Mr. or Ms. or Dr. or Prof. 
%
%
%  No changes in the next non-comment line also.
%
\begin{document}
%
%***************************************************
%
%  In the next line, replace "Seminar Title" by the title of your seminar.
%
\VAtitle{Seminar Title}%
%
%  In the next line, replace "Student" by your name. 
%  Write full name without, repeat without, Mr. or Ms.
%
\VAauthor{Student}%
%
%  In the next line, replace "2008" by the year of your admission 
%  to the College.
%
\VAadmissionyear{2010}%
%
%  In the next line, replace "VExxMCAnnn" by your 
%  University Examination Register Number.
%
\VAregisternumber{VExxxCSnnn}% 
%
%  In the next line, replace "Guide" by the full name of your guide or supervisor 
%  with Mr. or Ms. or Dr. or Prof.
%
\VAguide{Guide}% 
%
%  In the next line, replace "hod" by the full name of Head of Department 
%  with Mr. or Ms. or Dr. or Prof.
\VAhod{hod}
%
%  Printing the first few pages: title, certificate, acknowledgement, etc.
%  No changes in the next non-comment line.
%
%
\makeatletter
\renewcommand{\chapter}{\if@openright\cleardoublepage\else\clearpage\fi
                    %\thispagestyle{fancy}%
                    \global\@topnum\z@
                    \@afterindentfalse
                    \secdef\@chapter\@schapter}%
%
\makeatother
%
\pagestyle{empty}
%\frontmatter
%
%******************************************************
%
\begin{titlepage}
\newcommand{\HRule}{\rule{\linewidth}{0.5mm}}
\begin{center}
% Upper part of the page
% Title
%\quad\\[0.3cm]
{\large \sffamily CS 09 805 (P) : Seminar}\\[0.2cm]
\HRule \\[0.5cm]
{ \huge\sffamily \bfseries \vtitle}\\[0.2cm]
\HRule \\[5.25cm]
% Author
{\Large \sffamily\bfseries  \vauthor}\\
{\large \sffamily Reg. No. \vregisternumber\\ 
S8 CSE (\vadmissionyear\ Admissions)} \\[4.5cm] 
% Bottom of the page
%\vfill
\includegraphics[width=0.15\textwidth]{VidyaLogo}\\[0.3cm]
{\Large \sffamily\bfseries Department of Computer Science and Engineering}\\ {\large\sffamily Vidya Academy of Science \& Technology\\ \normalsize Thalakkottukara, Thrissur - 680 501}\\
({\tt http://www.vidyaacdemy.ac.in})\\
\end{center}
\end{titlepage}
%
%******************************************************
%

%
%******************************************************
%
\begin{titlepage}
\newcommand{\HRule}{\rule{\linewidth}{0.5mm}}
\begin{center}
% Upper part of the page
% Title
%\quad\\[0.3cm]
{\large \sffamily CS 09 805 (P) : Seminar}\\[1cm]
%\HRule \\[0.5cm]
{ \huge\sffamily \bfseries \vtitle}\\[5.9cm]
% Author
{\Large \sffamily \bfseries \vauthor}\\
{\large \sf  Reg. No. \vregisternumber\\ 
S8 CSE (\vadmissionyear\  Admissions)} \\[4cm] 
% Bottom of the page
%\vfill
\includegraphics[width=0.40\textwidth]{VidyaLogo}\\[0.2cm]
{\Large \sffamily \bfseries Department of Computer Science and Engineering}\\ {\large \sffamily Vidya Academy of Science \& Technology\\ \normalsize Thalakkottukara, Thrissur - 680 501}\\
({ \bf \tt http://www.vidyaacademy.ac.in})\\[1cm]
\end{center}
\end{titlepage}
%
%******************************************************
%
\clearpage
\pagenumbering{roman}
%
\addcontentsline{toc}{chapter}{\quad Certificate}

\begin{center}
\hrule
\quad\\[0.75cm]
{\Large \bf Department of Computer Science and Technology}\\
{\large \bf Vidya Academy of Science \& Technology}\\
{\normalsize \bf Thalakkottukara, Thrissur - 680 501\\
({\tt http://www.vidyaacademy.ac.in})}\\[0.5cm]
\includegraphics{VidyaLogo}\\[1.25cm]
%
{\Large \bf CERTIFICATE}\\[1cm]
%
\end{center}
\index{certificate}
\index{University of Calicut}
This is to certify that the report titled 
{\bf \vtitle} is a bona-fide record of the 
Seminar presented by {\bf \vauthor\  
(Reg. No. \vregisternumber)}    of the Eighth 
Semester B.Tech CSE (\vadmissionyear\ admissions) 
of Vidya Academy of Science \& Technology, 
Thrissur - 680 501 in partial fulfilment of the 
requirement for the award of the Bachelor of Technology in Computer Science and Engineering from the University of 
Calicut.\\[0.1cm]
%
\begin{center}
\begin{tabular}{llll}
\multicolumn{2}{l}{\bf Seminar Guide}    
&\multicolumn{2}{l}{\bf Head of Department}   \\
   & & \\
Name: &&  Name: \\
\vguide\ && \vhod\ \\ \\
Signature: \quad  ................................
\qquad \qquad && Signature: \quad 
................................ \\
& & & \\
Date: \hspace{10mm} \today & & Date: \hspace{10mm} \today \\
\end{tabular}
\\[1cm]
{\small (Seal of Department of Computer Science and Engineering)}
\quad \\[0.75cm]
\hrule
\end{center} 
%

%
\clearpage
\pagestyle{fancy}
\lhead{{\small \vtitle}} \rhead{\thepage}
\lfoot{{\scriptsize Department of Computer Applications, 
Vidya Academy of Science \& Technology, Thrissur}}
\rfoot{\includegraphics[width=0.5cm]{VidyaLogo}}
\cfoot{ }
%
\renewcommand{\headrulewidth}{1pt}
\renewcommand{\footrulewidth}{0.5pt}
%
\clearpage
\thispagestyle{fancy}
%
\addcontentsline{toc}{chapter}{\quad Acknowledgement}
\chapter*{Acknowledgement}
%
\index{acknowledgment}
I would like to thank the {\bf Lord Almighty}, the foundation of all wisdom who has been guiding me in every step. 

I wish to record my indebtedness and thankfulness to all who helped me to prepare this Seminar Report titled {\bf \vtitle } and present it in a satisfactory way. 


My sincere thanks to {\bf Dr.Sudha Balagopalan Ph.D, MISTE, MIEEE, MIET}, our Principal, for providing me all the necessary facilities.
I am also thankful to {\bf Col.M.P Induchudan(Retd) ME, FITE, FIE, SMCSI},the Head of Department,Computer Science and Engineering,
for his encouragement and valuable 
support.

I am especially thankful to my seminar guide {\bf  Ms.Geethu.P.C M.Tech, MISTE} and supervisor {\bf Ms.Livya George M.E} for giving me valuable suggestions and critical inputs in the preparation of the seminar report. 

My sincere thanks to my {\bf parents} for their valuable support and assitance for the work.My {\bf classmates} have always been 
helpful and I am grateful to them for 
patiently listening to  my presentation of the seminar and interacting with me through their valuable questions. 
%
\clearpage
%

%
%
%****************************************************
%
% No changes in the next two non-comment lines.
%
\chapter*{Abstract}
\addcontentsline{toc}{chapter}{\quad Abstract}
%
%  Here enter the abstract of the paper.
Bitcoin is the world's first decentralized digital currency, allowing the easy 
storage and transfer of cryptographic tokens, using a peer-to-peer network to carry
information, hashing as a synchronization signal to prevent double-spending, 
and a powerful scripting system to determine ownership of the tokens. There is a 
growing technology and business infrastructure supporting it.


A purely peer-to-peer version of electronic cash would allow online
payments to be sent directly from one party to another without going through a
financial institution. Digital signatures provide part of the solution, but the main
benefits are lost if a trusted third party is still required to prevent double-spending.

The network timestamps transactions by hashing them into an ongoing chain of
hash-based proof-of-work, forming a record that cannot be changed without 
redoing the proof-of-work. The longest chain not only serves as proof of the 
sequence of events witnessed, but proof that it came from the largest pool of CPU 
power. As long as a majority of CPU power is controlled by nodes that are not 
cooperating to attack the network, they'll generate the longest chain and outpace 
attackers. The network itself requires minimal structure. Messages are broadcast on 
a best effort basis, and nodes can leave and rejoin the network at will, accepting the 
longest proof-of-work chain as proof of what happened while they were gone.

To send money, you broadcast to the network that the amount on your account 
should go down, and the amount on a receiver’s account up. Nodes, or computers,
in the Bitcoin network apply that transaction to their copy of the ledger, and then 
pass on the transaction to other nodes. This, with some math-based security, is 
really all there is--a system that lets a group of computers maintain a ledger.

%
%***************************************************
%
%  The next line creates a Table of Contents. 
%  Do not delete this. There must be a Table of Contents.
%
\tableofcontents
%
%  The next line creates a page containing a List of Tables.
%  Delete it if there are no tables in your paper. 
%
\listoftables
%
%  The next line creates a page containing a List of Figures.
%  Delete it if there are no figures in your paper.
%
\listoffigures
%
%***************************************************
%
\mainmatter
%
%  The main contents of the paper begin here. 
%  In the following replace "Abcde Xyz" with 
%  the title of the first chapter of your paper.
%
\chapter{Introduction}
%
%  The text given below is only some filler text. Delete 
%  all those and replace them with the contents of your % 
%  paper. It is not exactly gibberish. (The filler text is a 
%  passage in Latin written by Cicero.)
%
Until Bitcoin’s invention in 2008 by the unidentified programmer known as Satoshi Nakamoto, online 
transactions always required a trusted third-party intermediary. For example, if Alice wanted to send \$100 
to Bob over the Internet, she would have had to rely on a third-party service like PayPal or MasterCard. 
Intermediaries like PayPal keep a ledger of account holders’ balances. When Alice sends Bob \$100, 
PayPal deducts the amount from her account and adds it to Bob’s account. Without such intermediaries, 
digital money could be spent twice. Imagine there are no intermediaries with ledgers, and digital cash 
is simply a computer file, just as digital documents are computer files. Alice could send \$100 to Bob 
by attaching a money file to a message. But just as with email, sending an attachment does not remove 
it from one’s computer. Alice would retain a copy of the money file after she had sent it. She could 
then easily send the same \$100 to Charlie. In computer science, this is known as the “double-spending” 
problem and until Bitcoin it could only be solved by employing a ledger-keeping trusted third party. Bit 
coin’s invention is revolutionary because for the first time the double-spending problem can be solved 
without the need for a third party. Bitcoin does this by distributing the necessary ledger among all the 
users of the system via a peer-to-peer network. Every transaction that occurs in the Bitcoin economy is 
registered in a public, distributed ledger, which is called the block chain. New transactions are checked 
against the block chain to ensure that the same bit coins haven’t been previously spent, thus eliminating 
the double-spending problem. The global peer-to-peer network, composed of thousands of users, takes the 
place of an intermediary; Alice and Bob can transact without PayPal. One thing to note right away is that 
transactions on the Bitcoin network are not denominated in dollars or Euros or yens they are on PayPal, 
but are instead denominated in Bitcoin. This makes it a virtual currency in addition to a decentralized 
payments network. The value of the currency is not derived from gold or government fiat, but from the 
value that people assign to it. The dollar value of a Bitcoin is determined on an open market, just as is the 
exchange rate between different world currencies.
%
%  Replace "Fghi Jklm" with the title of 
%  the first section of the first chapter of your Report.
%

\chapter{Related Works}
\textbf{Electronic funds transfer (EFT)} is the electronic exchange, transfer of money from one account to
another, either within a single financial institution or across multiple institutions, through computer-based 
systems.
It consists of different concepts;
\begin{itemize}
	\item Cardholder-initiated transactions, using a payment card such as a credit or debit card
	\item Direct deposit payment initiated by the payer
	\item Direct debit payments, sometimes called electronic checks, for which a business debits the consumer's bank accounts for payment for goods or services
	\item Electronic bill payment in online banking, which may be delivered by EFT or paper check
	\item Transactions involving stored value of electronic money, possibly in a private currency
	\item Electronic Benefit Transfer
\end{itemize}
EFTs include direct-debit transactions, wire transfers, direct deposits, ATM withdrawals and online 
bill pay services. Transactions are processed through the Automated Clearing House (ACH) network, 
the secure transfer system of the Federal Reserve that connects all U.S. banks, credit unions and other 
financial institutions.

For example, when you use your debit card to make a purchase at a store or online, the transaction 
is processed using an EFT system. The transaction is very similar to an ATM withdrawal, with near-
instantaneous payment to the merchant and deduction from your checking account.

Direct deposit is another form of an electronic funds transfer. In this case, funds from your employer’s
bank account are transferred electronically to your bank account, with no need for paper-based payment 
systems.

The increased use of EFTs for online bill payments, purchases and pay processes is leading to a paper-
free banking system, where a large number of invoices and payments take place over digital networks. 
EFT systems play a large role in this future, with fast, secure transactions guaranteeing a seamless transfer 
of funds within institutions or across banking networks.

EFT transactions, also known as an online transaction or PIN-debit transaction, also offer an alternative 
to signature debit transactions, which take place through one of the major credit card processing systems, 
such as Visa, MasterCard or Discover, and can cost as much as 3\% of the total purchase price. EFT 
processing, on the other hand, only charges an average of 1\% for debit card transactions.

\chapter{Existing System}
The first crypto currency to begin trading was Bitcoin in 2009. Since then, numerous crypto 
currencies have been created. Fundamentally, crypto currencies are specifications regarding the 
use of currency which seek to incorporate principles of cryptography to implement a distributed, 
decentralized and secure information economy. When comparing crypto currencies to fiat 
money, the most notable difference is in how no group or individual may accelerate, stunt or in 
any other way significantly abuse the production of money. Instead, only a certain amount of 
crypto currency is produced by the entire crypto currency system collectively, at a rate which is 
bounded by a value both prior defined and publicly known.

In centralized economic systems such as the Federal Reserve System governments regulate 
the value of currency by simply printing units of fiat money or demanding additions to digital 
banking ledgers, however governments cannot produce units of crypto currency and as such 
governments cannot provide backing for firms, banks or corporate entities which hold asset 
value measured in a decentralized crypto currency. The underlying technical system upon 
which all crypto currencies are now based was created by the anonymous group or individual 
known as Satoshi Nakamoto for the purpose of creating an economy within which the practice 
of fractional reserve banking would be fundamentally impossible.

Hundreds of crypto currency specifications now exist, most are similar to and derived from the 
first fully implemented crypto currency protocol, Bitcoin. Within crypto currency systems,
the safety, integrity, and balance of all ledgers is ensured by a swarm of mutually distrustful
parties, referred to as miners, who are, for the most part, general members of the public, actively
protecting the network by maintaining a high hash-rate difficulty for their chance at receiving
a randomly distributed small fee. Subverting the underlying security of a crypto currency is 
mathematically possible, but the cost may be unfeasibly high.

Most crypto currencies are designed to gradually introduce new units of currency, placing an 
ultimate cap on the total amount of currency that will ever be in circulation. This is done both to 
mimic the scarcity (and value) of precious metals and to avoid hyperinflation. As a result, such 
crypto currencies tend to experience hyper deflation as they grow in popularity and the amount 
of the currency in circulation approaches this finite cap. Compared with ordinary currencies 
held by financial institutions or kept as cash on hand, crypto currencies are less susceptible 
to seizure by law enforcement. Existing crypto currencies are all pseudonymous, though 
additions such as Zerocoin and its distributed laundry feature have been suggested, which would 
allow for anonymity.

\chapter{Proposed System}
Bitcoin is a form of digital currency, created and held electronically. No one controls it. Bitcoin aren’t
printed, like dollars or euros – they’re produced by lots of people running computers all around the world,
using software that solves mathematical problems. It’s the first example of a growing category of money 
known as cryptocurrency.

Bitcoin can be used to buy things electronically. In that sense, it’s like conventional dollars, euros, or yen, 
which are also traded digitally.

However, bit coin’s most important characteristic, and the thing that makes it different to conventional 
money, is that it is decentralized. No single institution controls the Bitcoin network. This puts some 
people at ease, because it means that a large bank can’t control their money.

This currency isn’t physically printed in the shadows by a central bank, unaccountable to the population, 
and making its own rules. Those banks can simply produce more money to cover the national debt, thus 
devaluing their currency.

Instead, Bitcoin is created digitally, by a community of people that anyone can join. Bitcoin are ‘mined’, 
using computing power in a distributed network. This network also processes transactions made with the 
virtual currency, effectively making Bitcoin its own payment network.

Bitcoin has several important features that set it apart from normal fiat currencies.

\begin{enumerate}
	\item It’s decentralized
	
	The Bitcoin network isn’t controlled by one central authority. Every machine that mines bitcoin and
	processes transactions makes up a part of the network, and the machines work together. That means that, 
	in theory, one central authority can’t tinker with monetary policy and cause a meltdown – or simply 
	decide to take people’s Bitcoin away from them, as the Central European Bank decided to do in Cyprus in 
	early 2013. And if some part of the network goes offline for some reason, the money keeps on flowing.
	\item It's easy to setup

	Conventional banks make you jump through hoops simply to open a bank account. Setting up merchant
	accounts for payment is another Kafkaesque task, beset by bureaucracy. However, you can set up a
	Bitcoin address in seconds, no questions asked, and with no fees payable.

	\item It's anonymous

	Well, kind of. Users can hold multiple Bitcoin addresses, and they aren’t linked to names, addresses, or
	other personally identifying information.

	\item It's completely transparent

	Bitcoin stores details of every single transaction that ever happened in the network in a huge version 
	of a general ledger, called the block chain. The block chain tells all. If you have a publicly used bitcoin
	address, anyone can tell how many Bitcoin are stored at that address. They just don’t know that it’s yours.
	There are measures that people can take to make their activities more opaque on the bitcoin network,
	though, such as not using the same Bitcoin addresses consistently, and not transferring lots of bitcoin to a
	single address.

	\item Transaction fees are miniscule

	Your bank may charge you a £10 fee for international transfers. Bitcoin doesn’t.

	\item It's fast

	You can send money anywhere and it will arrive minutes later, as soon as the bitcoin network processes
	the payment.

	\item It's non-repudiable

	When your bitcoins are sent, there’s no getting them back, unless the recipient returns them to you.
	They’re gone forever.
\end{enumerate}

\chapter{Bitcoin: A Peer-to-Peer Electronic Cash System}
\section{Introduction}

Commerce on the Internet has come to rely almost exclusively on financial institutions serving as
trusted third parties to process electronic payments. While the system works well enough for
most transactions, it still suffers from the inherent weaknesses of the trust based model.
Completely non-reversible transactions are not really possible, since financial institutions cannot
avoid mediating disputes. The cost of mediation increases transaction costs, limiting the
minimum practical transaction size and cutting off the possibility for small casual transactions,
and there is a broader cost in the loss of ability to make non-reversible payments for nonreversible
services. With the possibility of reversal, the need for trust spreads. Merchants must
be wary of their customers, hassling them for more information than they would otherwise need.
A certain percentage of fraud is accepted as unavoidable. These costs and payment uncertainties
can be avoided in person by using physical currency, but no mechanism exists to make payments
over a communications channel without a trusted party.

What is needed is an electronic payment system based on cryptographic proof instead of trust,
allowing any two willing parties to transact directly with each other without the need for a trusted
third party. Transactions that are computationally impractical to reverse would protect sellers
from fraud, and routine escrow mechanisms could easily be implemented to protect buyersA solution is proposed  to the double-spending problem using a peer-to-peer distributed
timestamp server to generate computational proof of the chronological order of transactions. The
system is secure as long as honest nodes collectively control more CPU power than any
cooperating group of attacker nodes.

\section{Transactions}

We define an electronic coin as a chain of digital signatures. Each owner transfers the coin to the
next by digitally signing a hash of the previous transaction and the public key of the next owner
and adding these to the end of the coin. A payee can verify the signatures to verify the chain of
ownership.

\begin{figure}[ht!]
\centering
\includegraphics[trim = 0mm 0mm 30mm 0mm, width=120mm]{transaction}
\caption{Transaction}
\end{figure}

The problem of course is the payee can't verify that one of the owners did not double-spend
the coin. A common solution is to introduce a trusted central authority, or mint, that checks every
transaction for double spending. After each transaction, the coin must be returned to the mint to
issue a new coin, and only coins issued directly from the mint are trusted not to be double-spent.
The problem with this solution is that the fate of the entire money system depends on the
company running the mint, with every transaction having to go through them, just like a bank.
We need a way for the payee to know that the previous owners did not sign any earlier
transactions. For our purposes, the earliest transaction is the one that counts, so we don't care
about later attempts to double-spend. The only way to confirm the absence of a transaction is to
be aware of all transactions. In the mint based model, the mint was aware of all transactions and
decided which arrived first. To accomplish this without a trusted party, transactions must be
publicly announced [1], and we need a system for participants to agree on a single history of the
order in which they were received. The payee needs proof that at the time of each transaction, the
majority of nodes agreed it was the first received.

\section{Timestamp Server}

The solution we propose begins with a timestamp server. A timestamp server works by taking a hash of a block of items to be timestamped and widely publishing the hash, such as in a newspaper or Usenet post. The timestamp proves that the data must have existed at the time, obviously, in order to get into the hash. Each timestamp includes the previous timestamp in its hash, forming a chain, with each additional timestamp reinforcing the ones before it.

\begin{figure}[ht!]
\centering
\includegraphics[trim = 0mm 0mm 0mm 0mm, width=120mm]{timestamp}
\caption{Timestamp Chain}
\end{figure}

\section{Proof-of-Work}

To implement a distributed timestamp server on a peer-to-peer basis, we will need to use a proof-of-work system similar to Adam Back's Hashcash, rather than newspaper or Usenet posts. The proof-of-work involves scanning for a value that when hashed, such as with SHA-256, the hash begins with a number of zero bits. The average work required is exponential in the number of zero bits required and can be verified by executing a single hash.
For our timestamp network, we implement the proof-of-work by incrementing a nonce in the block until a value is found that gives the block's hash the required zero bits. Once the CPU effort has been expended to make it satisfy the proof-of-work, the block cannot be changed without redoing the work. As later blocks are chained after it, the work to change the block would include redoing all the blocks after it.

\begin{figure}[ht!]
\centering
\includegraphics[trim = 0mm 0mm 0mm 0mm, width=120mm]{block_chain}
\caption{Block Chain}
\end{figure}

The proof-of-work also solves the problem of determining representation in majority decision making. If the majority were based on one-IP-address-one-vote, it could be subverted by anyone able to allocate many IPs. Proof-of-work is essentially one-CPU-one-vote. The majority decision is represented by the longest chain, which has the greatest proof-of-work effort invested in it. If a majority of CPU power is controlled by honest nodes, the honest chain will grow the fastest and outpace any competing chains. To modify a past block, an attacker would have to redo the proof-of-work of the block and all blocks after it and then catch up with and surpass the work of the honest nodes. We will show later that the probability of a slower attacker catching up diminishes exponentially as subsequent blocks are added.

To compensate for increasing hardware speed and varying interest in running nodes over time, the proof-of-work difficulty is determined by a moving average targeting an average number of blocks per hour. If they're generated too fast, the difficulty increases.

\section{Network}

The steps to run the network are as follows:
\begin{enumerate}
	\item New transactions are broadcast to all nodes.
	\item Each node collects new transactions into a block.
	\item Each node works on finding a difficult proof-of-work for its block
	\item When a node finds a proof-of-work, it broadcasts the block to all nodes.
	\item Nodes accept the block only if all transactions in it are valid and not already spent.
	\item Nodes express their acceptance of the block by working on creating the next block in the chain, using the hash of the accepted block as the previous hash.
\end{enumerate}

Nodes always consider the longest chain to be the correct one and will keep working on extending it. If two nodes broadcast different versions of the next block simultaneously, some nodes may receive one or the other first. In that case, they work on the first one they received, but save the other branch in case it becomes longer. The tie will be broken when the next proof- of-work is found and one branch becomes longer; the nodes that were working on the other branch will then switch to the longer one.

New transaction broadcasts do not necessarily need to reach all nodes. As long as they reach many nodes, they will get into a block before long. Block broadcasts are also tolerant of dropped messages. If a node does not receive a block, it will request it when it receives the next block and realizes it missed one.

\section{Incentive}

By convention, the first transaction in a block is a special transaction that starts a new coin owned by the creator of the block. This adds an incentive for nodes to support the network, and provides a way to initially distribute coins into circulation, since there is no central authority to issue them. The steady addition of a constant of amount of new coins is analogous to gold miners expending resources to add gold to circulation. In our case, it is CPU time and electricity that is expended.

The incentive can also be funded with transaction fees. If the output value of a transaction is less than its input value, the difference is a transaction fee that is added to the incentive value of the block containing the transaction. Once a predetermined number of coins have entered circulation, the incentive can transition entirely to transaction fees and be completely inflation free.

The incentive may help encourage nodes to stay honest. If a greedy attacker is able to assemble more CPU power than all the honest nodes, he would have to choose between using it to defraud people by stealing back his payments, or using it to generate new coins. He ought to find it more profitable to play by the rules, such rules that favour him with more new coins than everyone else combined, than to undermine the system and the validity of his own wealth.

\section{Reclaiming Disk Space}

Once the latest transaction in a coin is buried under enough blocks, the spent transactions before it can be discarded to save disk space. To facilitate this without breaking the block's hash, transactions are hashed in a Merkle Tree, with only the root included in the block's hash. Old blocks can then be compacted by stubbing off branches of the tree. The interior hashes do not need to be stored.

\begin{figure}[ht!]
\centering
\includegraphics[trim = 0mm 0mm 0mm 0mm, width=120mm]{merkle_tree_and_pruning}
\caption{Transactions hashed in a Merkle Tree and pruning}
\end{figure}

A block header with no transactions would be about 80 bytes. If we suppose blocks are generated every 10 minutes, 80 bytes * 6 * 24 * 365 = 4.2MB per year. With computer systems typically selling with 2GB of RAM as of 2008, and Moore's Law predicting current growth of 1.2GB per year, storage should not be a problem even if the block headers must be kept in memory.

\section{Simplified Payment Verification}

It is possible to verify payments without running a full network node. A user only needs to keep a copy of the block headers of the longest proof-of-work chain, which he can get by querying network nodes until he's convinced he has the longest chain, and obtain the Merkle branch linking the transaction to the block it's timestamped in. He can't check the transaction for himself, but by linking it to a place in the chain, he can see that a network node has accepted it, and blocks added after it further confirm the network has accepted it.

\begin{figure}[ht!]
\centering
\includegraphics[trim = 0mm 0mm 0mm 0mm, width=120mm]{proof_of_work}
\caption{Proof-of-Work Chain}
\end{figure}

As such, the verification is reliable as long as honest nodes control the network, but is more vulnerable if the network is overpowered by an attacker. While network nodes can verify transactions for themselves, the simplified method can be fooled by an attacker's fabricated transactions for as long as the attacker can continue to overpower the network. One strategy to protect against this would be to accept alerts from network nodes when they detect an invalid block, prompting the user's software to download the full block and alerted transactions to confirm the inconsistency. Businesses that receive frequent payments will probably still want to run their own nodes for more independent security and quicker verification.

\section{Combining and Splitting Value}

Although it would be possible to handle coins individually, it would be unwieldy to make a separate transaction for every cent in a transfer. To allow value to be split and combined,transactions contain multiple inputs and outputs. Normally there will be either a single input from a larger previous transaction or multiple inputs combining smaller amounts, and at most two outputs: one for the payment, and one returning the change, if any, back to the sender.

\begin{figure}[ht!]
\centering
\includegraphics[trim = 0mm 0mm 0mm 0mm, width=70mm]{multiple_in_outs_transaction}
\caption{Inputs and Outputs in a transaction}
\end{figure}

It should be noted that fan-out, where a transaction depends on several transactions, and those transactions depend on many more, is not a problem here. There is never the need to extract a complete standalone copy of a transaction's history.

\section{Privacy}

The traditional banking model achieves a level of privacy by limiting access to information to the parties involved and the trusted third party. The necessity to announce all transactions publicly precludes this method, but privacy can still be maintained by breaking the flow of information in another place: by keeping public keys anonymous. The public can see that someone is sending an amount to someone else, but without information linking the transaction to anyone. This is similar to the level of information released by stock exchanges, where the time and size of individual trades, the \"tape\", is made public, but without telling who the parties were.

\begin{figure}[ht!]
\centering
\includegraphics[trim = 0mm 0mm 0mm 0mm, width=120mm]{traditional_privacy_model}
\caption{Privacy Models}
\end{figure}

As an additional firewall, a new key pair should be used for each transaction to keep them from being linked to a common owner. Some linking is still unavoidable with multi-input transactions, which necessarily reveal that their inputs were owned by the same owner. The risk is that if the owner of a key is revealed, linking could reveal other transactions that belonged to the same owner.


%
% Continue adding new chapters, sections and 
% subsections.
%
% ***   Creating Bibliography   ***
%
% Do not delete the following two lines
%
\clearpage
\addcontentsline{toc}{chapter}{\quad Bibliography}
%
% The markups for creating the bibliography begins here.
% One sample bibliography item is included. 
% Add additional items if there are any.
%
\begin{thebibliography}{99}
%
\bibitem{First}
{\em First Lessons in \LaTeX}, by Dr. V N Krishnachandran, 
Vidya Academy of Science \& Technology, 
Thrissur - 680 501, 2011.
%
\end{thebibliography}
%
% Printing the last page
%
%
\newpage
\thispagestyle{empty}
\vspace*{\fill}
\begin{flushright}
\includegraphics{VidyaLogo}\\[0.5cm]
{\Large \bf \sf  Department of Computer Science and Engineering}\\
{\sf Vidya Academy of Science \& Technology\\
Thalakkottukara, Thrissur - 680 501\\
({\tt http://www.vidyaacademy.ac.in})}
\end{flushright}
%
%
%   ***   The end   ***
%
\end{document}
